\section{ĐỀ SỐ \thesection } %ĐỀ 4
\setcounter{bt}{0}

%%Câu 1.
\begin{bt} Thực hiện phép tính:
	\begin{listEX}[2]
		\item $(-42)+27+(-10)$;
		\item $5^{6}: 5^{4}+2^{3} \cdot 2^{2}-1^{2017}$.
	\end{listEX}
	
	\loigiai{a) -25 \hfill b) 56
	}
\end{bt}


%%Câu 2.
\begin{bt} Tìm số nguyên $x$, biết:\begin{listEX}[2]
		\item $22+x=3$;
		\item $48-3(x+5)=24$.
	\end{listEX}
	
	\loigiai{a) -19 \hfill b) 3}
\end{bt}


%%Câu 3.
\begin{bt} Số đội viên của một liên đội khoảng từ $300$ đến $400$ đội viên. Mỗi lần xếp hàng $12$, hàng $15$, hàng $18$ đều vừa đủ không thừa ai. Hỏi liên đội đó có bao nhiêu đội viên?
	
	\loigiai{360 đội viên}
\end{bt}


%%Câu 4.
\begin{bt} 
	\begin{multicols}{2}
		Biểu đồ tranh bên cạnh cho biết số lượt ô tô vào gửi tại một bãi đỗ xe vào các ngày trong một tuần.
		\begin{listEX}
			\item Tính số xe được gửi tại bãi đỗ xe trong ngày thứ Tư.
			\item Tính tổng số xe đã gửi tại bãi từ thứ Ba đến ngày thứ Năm trong một tuần.
		\end{listEX}
		\begin{center}
			\def\icon{\twemoji[scale=0.8]{car}}
			\def\vetranh#1{\foreach \dem in{1,...,#1}{\icon \ }}{
				\begin{tabular}{|c|l|}\hline 
					Thứ Hai &\vetranh{5}\\ \hline 
					Thứ Ba &\vetranh{7} \\ \hline 
					Thứ Tư &\vetranh{3} \\ \hline 
					Thứ Năm &\vetranh{4} \\ \hline 
					Thứ Sáu&\vetranh{6} \\ \hline 
				\end{tabular}\\
				(\icon = 3 xe ô tô)}
		\end{center}
	\end{multicols}
	\loigiai{a) $9$ ô tô \hfill b) $42$ ô tô
	}
\end{bt}

%%Câu 5.
\begin{bt} 
	\immini[thm]{
		Một mảnh vườn hình chữ nhật có chiều dài $10$ m; chiều rộng bằng $7$ m. Người ta làm $2$ lối đi rộng $1$ m như hình vẽ. Phần đất còn lại dùng để trồng cây.
		\begin{listEX}
			\item Tính diện tích mảnh vườn hình chữ nhật.
			\item Tính diện tích phần đất dùng để trồng cây.  
		\end{listEX}
	}{
		\begin{tikzpicture}[scale=1,line width=0.8pt,line join=round,>=stealth]
			\def\a{4.2}  \def\b{6}  \def\c{0.6}  \def\d{0.3}
			\path (0,0) coordinate (A)
			(\b,0) coordinate (B)
			(\b,\a) coordinate (C)
			(0,\a) coordinate (D)
			;
			\fill[green!30!gray] (A) rectangle (C);
			\fill[white] 
			(0,0.5*\a-0.5*\c) rectangle (\b,0.5*\a+0.5*\c)
			(0.5*\b-0.5*\c,0) rectangle (0.5*\b+0.5*\c,\a);
			\draw        (\c,\a/2) node{Lối đi}
			(\b/2,\a-\c) node[rotate=-90]{Lối đi}
			(\b/2,\c) node[rotate=-90]{Lối đi}
			;
			\draw[latex-latex] (A)++(-90:\d)--++(0:\b) node[midway,below]{10 m};
			\draw[latex-latex] (\b/2-\c/2,\a+\d)--++(0:\c) node[midway,above]{1m};
			\draw[latex-latex] (A)++(180:\d)--++(90:\a/2-\c/2) node[midway,left]{3m};
			\draw[latex-latex] (D)++(180:\d)--++(-90:\a/2-\c/2) node[midway,left]{3m};
			\draw[latex-latex] (\b+\d,\a/2-\c/2)--++(90:\c) node[midway,right]{1m};
		\end{tikzpicture}
	}
	
	\loigiai{a) $70 \text{ m}^2$ \hfill b) $54 \text{ m}^2$
	}
\end{bt}



%%Câu 6.
\begin{bt} Tìm số tự nhiên $n$ để biểu thức $A=\frac{15}{2 n+1}$ có giá trị là một số tự nhiên.
	
	\loigiai{$n\in\{0, 1, 2, 7\}$}
\end{bt}
